The physical findings summarised in the preceding sections rest on a mature
toolkit for simulating, measuring, and locating drainage systems. Across
site-scale validations, benchmark comparisons of analytical solutions,
pipe-scale experiments, watershed representations, formal calibration and
uncertainty analyses, and geospatial detection, a robust set of consolidated
findings has emerged: process models reproduce drainage flow and nitrate
reliably after calibration but with skill that improves strongly with temporal
aggregation; treating drainage as purely saturated flow over-predicts
discharge; saturated conductivity and evapotranspiration dominate model
sensitivity; data-driven and remote-sensing methods now complement process
models; and no single geophysical or imaging technique is universally best for
mapping buried networks. This section consolidates that methodological
evidence so the physical and management conclusions can stand without model
mechanics.

\subsection{Field-scale drainage and water-quality models}

\textsc{DRAINMOD} is by far the most extensively tested field-scale
water-management model and is repeatedly judged a reliable design and
evaluation tool for poorly and artificially drained lands, its process-based
structure with independently measurable inputs describing drainage hydrology
across many climates when soil hydraulic properties are adequately
characterised \parencite{Singh2019b,Gowda2012}. A defining, cross-cutting
result is that skill is scale-dependent in time: water-table depth and
drainage are reproduced with accuracy improving markedly from daily to monthly
to annual scale. In North Carolina, water-table and daily-drainage
Nash--Sutcliffe efficiencies were 0.49--0.73, monthly drainage 0.77--0.87, and
annual volumes within 5\% of observations \parencite{Gowda2012}; Iowa and
eastern-Canada cumulative drainage fell within 2--10\% and about 30~mm of
measurements \parencite{Singh2006b,Dayyani2009}. The persistent weak point is
daily peak flow, which is systematically underestimated
\parencite{Morrison2014,Singh1996}, and uncertainty widths grow from daily to
yearly scale while annual predictions remain within bounds \parencite{Wang2006}.
The model has been validated well beyond its humid origin---cold Canadian
snowmelt ($R^2$ up to 0.93), semi-arid India, and Illinois claypan soils
\parencite{Dayyani2009,GUPTA1993,MOSTAGHIMI1989,Abduljaleel2023,Awad2022b}---but
delineated failure modes include unaccounted deep and lateral seepage
\parencite{Sinai2006}, absent time-variable depression storage
\parencite{Abduljaleel2023}, and its parallel-drain geometry requiring
plot-by-plot calibration in irregular landscapes \parencite{He2002}. Where
site data are scarce, pedotransfer estimation (ROSETTA, van Genuchten fits)
with automated calibration (PEST) drives it reliably
\parencite{Singh2006b,Qi2015,Malota2024,Malota2015}, and enhancements target
the daily-peak weakness through macropore flow (capturing $\sim$75\% of peak
events) and neural-network surrogates \parencite{Askar2020,Awad2022,Awad2022b};
the evapotranspiration algorithm is a pronounced sensitivity, the best ET
method not necessarily best for water-table prediction
\parencite{Malakshahi2020}.

Richards-equation and comparative codes generally perform comparably.
HYDRUS~(2D/3D) predicts fluxes and water tables closely (EF 0.91--1) and
clarifies mechanisms such as conductivity control of drawdown in compacted
soils \parencite{Karandish2018,DarziNaftchali2018,Mante2018,Guedessou2021,Ebrahimian2015};
CSUID, WAVE, SWAP, \textsc{WaSim}, CATHY, RZWQM, WEPP, ADAPT, and the
few-parameter \textsc{SIDRA} family have all calibrated satisfactorily, with
\textsc{DRAINMOD} often giving the most accurate discharge among them
\parencite{Alzraiee2013b,Alzraiee2013,ElSadek2001,Prasher1996,Verma2014,Malota2022,Muma2017,Malone2014,Xian2017,RevueltaAcosta2021,Mendez2004,Zimmer2023,Henine2022,LESAFFRE1988}.

For water quality the same pattern recurs across the model families
(\textsc{DRAINMOD} suite, RZWQM/RZWQM2, SWAT/APEX, APSIM, DNDC, HYDRUS): after
calibration they reproduce annual and cumulative drainage and nitrate loss
well---frequently within $\pm5$--$20\%$---and capture the direction of
management effects, but predictive skill falls with temporal disaggregation,
nitrate is predicted less accurately than flow, and daily nitrate
concentrations, crop N uptake, and residual soil N are recurrent weak points
\parencite{Youssef2018,Negm2014,Singh2020b,Dayyani2010,Ma2001,Liang2018,Ikenberry2017,Marjerison2016,Ojeda2018,Qi2025,Ma2007,Ahuja2007}.
\textsc{DRAINMOD-N~II} and RZWQM2 both satisfactorily predict tile nitrate,
and \textsc{DRAINMOD-DSSAT} predicts drainage, nitrate, and yield very well
\parencite{Du2017,Qi2012,Negm2017,Singh2022c}; cropping-system models
(DSSAT, RZWQM, DNDC, APSIM, SWAT) likewise reproduce cover-crop and
rotation effects on nitrate loss \parencite{Jeong2018,Ahmed2007b,Ale2012b,Pasley2021,Motarjemi2021}.
Phosphorus extensions are advancing but remain limited: RZWQM2-P first
simulated daily dissolved reactive P and its pathway partitioning, while
\textsc{DRAINMOD-P} predicted discharge and total P but poorly predicted DRP,
and few standard models yet represent tile P
\parencite{Sadhukhan2019,Grewal2025,BinShokrana2023b,Askar2021,Askar2025,Radcliffe2015,Wang2020b}.
Cold-season snowmelt and freeze--thaw errors motivated dedicated developments
(\textsc{DRAINMOD}~5.1, DRAIN-WARMF, RZ-SHAW), and calibration studies
repeatedly identify lateral saturated conductivity in the tile layer and field
capacity as controlling predictions
\parencite{Wang2006c,Dayyani2010b,Li2024b,Ma2007c,Malone2017}.

\subsection{Flow physics: saturated and unsaturated flow and preferential-flow representation}

A central mechanistic result is that treating drainage as purely saturated
flow overstates discharge. A saturated-flow model over-predicted seepage and
drainage time by roughly an order of magnitude relative to experiment, an
analytical solution, and the SUTRA Richards-equation model, because a vadose
zone retaining substantial water persists even at steady state
\parencite{Xin2016b}. Saturated--unsaturated drainage is governed by
gravitational, pressure, capillary, and air-pressure potentials, so saturated
conductivity alone is an inadequate design basis \parencite{Li2021}. Field
diagnosis in Japanese paddies showed that vertical, not horizontal, resistance
limits the maximum drainage rate \parencite{Nishida2020}, accounting for
vertical head loss improved water-table prediction \parencite{Talukolaee2022},
and retaining unsaturated-zone buffering matched numerical models closely
whereas ignoring it distorted transients \parencite{VandeCraats2021};
equivalent-medium representations reproduce integrated outflow well but can
misrepresent internal dynamics where fast vertical flow occurs
\parencite{Carlier2007}. At field-to-hillslope scale the representation of
macropores is decisive for reproducing rapid drainage. Dual-permeability and
dual-porosity models reproduce hydrographs well ($R^2=0.84$; groundwater
$r=0.98$) while single-porosity models reach acceptable accuracy with far
fewer parameters, though all fail to match low dry-season discharge
\parencite{Varvaris2018}; dual-permeability formulations enable low matrix
conductivity while still reproducing fast events
\parencite{Warsta2014,Salo2017,Koivusalo1999}, adding macropore flow improves
peak simulation without changing total drainage \parencite{Zhou2013}, and
bimodal retention functions predict preferential flow far better than unimodal
ones \parencite{Mohanty1997}. Fully coupled surface--subsurface--tile models
are recommended for detailed dynamics, with equivalent-porous-medium or
seepage-node simplifications reproducing outflow almost exactly for watershed
integration \parencite{DeSchepper2015,Boico2022,Xian2017}.

\subsection{Analytical drain-spacing solutions and pipe-scale hydraulics}

Unsteady-state (transient) spacing equations yield markedly narrower spacings
than steady-state equations and are preferred under variable rainfall---about
30~m versus 55~m in Iranian paddies---with the appropriate transient model
depending on rainfall duration and design criterion
\parencite{OOSTERBAAN1988,DarziNaftchally2014,Moshayedi2020,Pali2014,Nédélec2010}.
Incorporating evapotranspiration and unsaturated-zone flow generally permits
wider, more economical spacing: ET-aware solutions allow 9--18\% increases,
bare-soil evaporation 23--74\%, and unsaturated horizontal flow up to 30\% for
coarse soils near a shallow impervious layer
\parencite{NIKAM1992,Singh2010,Upadhyaya2000,AbdelFattah2005,Yousfi2014}.
Bi-level and sloping-aquifer solutions show asymmetric water tables and
slightly to moderately larger spacings, and dual-depth systems always remove
water more effectively than single-depth ones
\parencite{Verma1998,Kankarej2023,Upadhyaya2013,Hartani2001,Kirkham1997}.
Analytical solutions for multi-layered and ponded systems have been verified
against experiments and MODFLOW, establishing strong three-dimensionality and
sensitivity to anisotropy near drains, while flat-land theory remains
applicable on slopes up to about $0.15$~m~m$^{-1}$
\parencite{Chakrabarti2020,Sarmah2015,SHARMA1991,Afruzi2014,Lin2025,CHU1977,FIPPS1989,LORRE1994};
no single linear-diffusion model, however, can describe all stages of drawdown
\parencite{Kacimov2000}. The most complete descriptions solve the Richards
equation, but corrected Boussinesq approximations remain useful for design when
a variable drainable porosity or effective specific yield is used
\parencite{TANG1977,Xin2016c,VANDELEUR1959,SINGH1980,Bouarfa2000}. Single-tile
solutions using a variable region of influence eliminate the spacing parameter
for irregular systems \parencite{Cooke2001}, and space-fractional and inverse
formulations improve fits in heterogeneous soils and circumvent the awkward
drainable-porosity parameter
\parencite{FathAllah2025,Mehdinejadiani2013,Mehdinejadiani2020,Mehdinejadiani2022,Samani2007,Chahar2010,SINGH1992,KUMAR1994,Onyejekwe1999,Bahceci2008}.

At the pipe scale, entrance resistance and effective radius are now quantified
across perforation and envelope geometries: 3--4-row pipes averaged 0.6~cm
effective radius and 8-row reached 1.9~cm, while a knitted sock raised it to
5.7--6.0~cm and nearly eliminated entrance head loss, increasing perforation
rows being the most efficient improvement
\parencite{Ghane2022b,Ghane2022}. Numerical and experimental work shows
rectangular slots are more hydraulically advantageous and structurally safer
than circular holes, and screen envelopes reduce resistance and prevent rapid
clogging
\parencite{Gaj2020,Oyarce2017,Gaj2021,Stuyt2006}. Within-conduit flow is a
spatially varied flow problem in which Manning's equation alone is inadequate,
local head loss is non-negligible, and resistance rises 40--50\% over
gradually varied flow
\parencite{Bayat2011,Graber2004,Kouchakzadeh2011,Nédélec2008,Wu2022,Gupta2019}.

\subsection{Watershed representation, sensitivity, calibration, uncertainty, and evaluation}

How drains and subsurface heterogeneity are conceptualised strongly controls
simulated catchment behaviour: two grid conceptualisations of a Danish
catchment produced drain contributions of 13\% versus 31\% despite both
matching total streamflow within 0.1\% \parencite{Karan2021}. Seepage-node
representations are faster and better suited to unknown drain locations
\parencite{Boico2022}, distributed HydroGeoSphere reproduced flashy heads
better than MIKE SHE \parencite{DeSchepper2017}, machine learning links
streamflow signatures to simulated drain fraction \parencite{Schneider2022},
embedding drainage in integrated models reproduced streamflow without
site-specific calibration \parencite{Rathore2024,Muma2016}, and equivalent-medium
representations match explicit tile simulations within 2\%
\parencite{Thomas2016}; standard SWAT under-represents subsurface drainage
without explicit adaptation \parencite{Manguerra1998,Golmohammadi2016,Rahman2014}.
For regional assessment, MODFLOW matched observed heads closely, and new
Hooghoudt-based subsurface-drainage packages markedly outperform the
conventional DRN package (RRMSE $<5\%$) by calibrating only conductivity
\parencite{Mirlas2012,Ghorbani2016,Yang2022b,Han2025,Qian2025b}.

Formal analyses converge on saturated hydraulic conductivity and
evapotranspiration estimation as the dominant sources of sensitivity and
uncertainty, with van Genuchten shape parameters, drain spacing, and depth
also prominent, and parameter uncertainty propagating into predicted salinity
and yield
\parencite{NabiSichani2012,Malakshahi2020,Wang2006,Sedaghatdoost2018,Timalsina2023,Alzraiee2013}.
Objective-function design strongly affects identifiability: genetic algorithms
generally outperform PEST, combining discharge with water-table depth already
gives excellent results, and variational data assimilation cut a five-year
cumulative-volume discrepancy below 1\%
\parencite{Sedaghatdoost2019,Chelil2022b,Schoups2005}. The review literature
concludes no universal model exists and that automatic calibration is
preferable because it supports uncertainty estimation \parencite{Singh2019b}.
Crucially, evaluation choices are themselves substantive: the Kling--Gupta
efficiency consistently outperformed and could reverse rankings relative to
Nash--Sutcliffe, and regression-based tests exposed gross over-prediction
invisible to $R^2$
\parencite{Bruhn2025,STONE1995,SAVABI1995,Wessel2022}. Measurement protocols
matter equally: calibrated V-notch and compound weirs, flow meters, and
flow-proportional sampling minimise bias in discharge and load estimates that
otherwise propagates into nutrient budgets
\parencite{Christianson2019,BinShokrana2021,Napierala2023,Katuwal2024,Enamorado2007,Dialameh2022,Young2023}.

\subsection{Data-driven, remote-sensing, and geostatistical methods}

Data-driven methods frequently match or outperform process models for the
specific outputs on which they are trained while requiring less data and
computation. Neural networks predicted real-time water-table control and
saturated conductivity more accurately than pedotransfer functions
(the ANN prediction of saturated conductivity reaching $R^2$ up to 0.96),
random forests outperformed linear regression
for agronomic prediction, and an Attention~U-Net surrogate reached test $R^2$
of 0.997--1 for head, drainage volume, and salt discharge while drastically
reducing computation
\parencite{Yang2000,Ellafi2021,Djiemon2019,Lei2026,Sarangi2006,Bakhsh2005b,Bakhsh2007b}.
Parsimonious surrogates (Cubist, SVM, NIT-DRAIN) can match process models for
flow while training far faster, though process models remain preferable for
climate extrapolation
\parencite{Sharma2003,Negm2016,Frederiksen2023,Chelil2022,Li2025}. Remote
sensing combined with deep learning now maps drained extent at regional to
national scales with accuracies commonly exceeding 85--96\% and totals within a
few percent of census estimates: AgTile-US reported 22.34~Mha ($-0.35\%$ versus
USDA) and SEETileDrain achieved 0.96 testing accuracy
\parencite{Valayamkunnath2020,Wan2024,Moller2018}, while convolutional
segmentation networks reliably delineate visible drain lines and outperform
classical edge detectors \parencite{Redoloza2023,Song2021,Woo2023,Naz2009}.
Ancillary sensing infers drainage-relevant surface conditions through SAR,
Landsat SWIR, and vegetation indices
\parencite{Meunier2001,Rijal2013,Xu2018,Khanal2020}. Finally, kriging-family
methods are consistently the most reliable interpolators for the spatially
variable soil properties governing drain spacing: in the Nile Delta,
conductivity was isotropic with a 5.6~km range and kriged values were
recommended as design inputs, while accounting for variability widened spacing
about 5\% and saved substantial pipe, and neglecting $K_{sat}$ variability
could bias design spacing by $-27\%$ to $+3\%$
\parencite{GALLICHAND1991,GALLICHAND1992,GALLICHAND1993,Moustafa1998b,Moustafa2000,Moustafa1998,Akis2015,Sepaskhah2004}.

\subsection{Mapping drainage infrastructure and quantifying extent}

Because buried networks are frequently undocumented, a substantial body of work
establishes that non-invasive remote sensing and near-surface geophysics can
detect and map subsurface drainage, with a consistent conclusion that no single
method or spectral band is universally superior and that detection depends
strongly on soil moisture, soil composition, drain depth and diameter, and
survey timing relative to rainfall and field operations
\parencite{Carlsen2024,Wienken2024,Koganti2021}. In the largest multi-site
assessment (29 sites), UAV visible-color imagery detected drain lines at 48\%
of sites, multispectral at 59\%, and thermal infrared (TIR) at 69\%, so
acquiring all three bands is best practice \parencite{Allred2020}. TIR is the
most consistently effective single band, detecting roughly 60\% of
infrastructure even under dry conditions when visible imagery failed, though it
exploits apparent temperature differences below 1.0\,\textdegree C, requires
careful timing (optimally within 12~h of $\sim$2.5~mm rain and late morning to
early afternoon), and can fail under low contrast, where RGB imagery guides
interpretation
\parencite{Allred2018,Woo2019,Allred2021,Allred2022,Kratt2020}. Aerial visible
and multispectral photography likewise reveals functional drains through soil
and vegetation contrasts and has repeatedly exposed systems absent from
official registries \parencite{Tlapáková2017}. Ground-penetrating radar senses
pipes directly, with 250--1000~MHz antennas effective, antenna orientation
interacting strongly with soil moisture, and RTK/GNSS-positioned surveys
resolving system type in only four to six hours without dense grids; 3D-GPR
succeeds where soil electrical conductivity is low enough to limit attenuation
\parencite{Allred2005,Allred2013,Allred2018b,Koganti2020}. Geoelectrical
methods coupling time-lapse ERT with injected water or pressurised air locate
networks by induced resistivity change
\parencite{Jouen2018,Henine2021}. The best-supported strategy is combining
methods with complementary sensitivities---particularly UAV RGB/TIR imagery
with GPR, which adds pipe depth and orientation---since each has distinct
favourable conditions and linear farm features are otherwise easily confused
with drain lines \parencite{Koganti2021,Carlsen2024,Wienken2024,Allred2020}.

Beyond field scale, machine learning and hydrologic filtering extend mapping to
watersheds and regions: satellite big data with random forests map potential
drainage, and combining remote sensing, decision-tree classification, and
groundwater-table information improves watershed delineation
\parencite{Cho2019,Yimer2023}. Once layouts are known, GIS/GPS and
DEM-based tools support design and management; a suite of QGIS plugins tested in
Illinois streamlined identifying landscape features, producing topologically
sound flow networks, and optimising pipe depth and size in under ten minutes
versus over two hours for proprietary software
\parencite{Shedekar2018,Falasy2025,Falasy2024,Falasy2025b}. Quantifying extent
frames the significance of drainage effects: regional mapping estimated about
25\% of Germany drained (30--40\% on hydromorphic soils), U.S. datasets diverge
at watershed scale so best estimates remain locally conditioned, and this
extent rests on validated salt- and water-balance models, empirical
effluent-quality relationships, and geophysical (EMI/apparent-EC) monitoring
that together allow design principles to transfer with confidence to new sites
\parencite{Hirt2005,Hirt2011,Jame2022,Chen2025,Tian2023,Bahçeci2006,Skaggs2006,Narjary2024,Qin2025}.